\section{Implementation}
\label{sec:impl}

The RC1 artifact is a Rust workspace plus an independent Linux kernel
implementation. Tagged source and release artifacts are available from the
project repository.  Published Cargo packages use the \texttt{zero-sfs-*}
prefix~\citep{SfsArtifact2026}. \Cref{tab:artifact} summarizes the paths
relevant to this paper. Operational service features such as proxy deployment,
rate limiting, health endpoints, and recovery-code handling are implemented
but are not treated as research contributions.

\begin{table}[t]
  \centering
  \small
  \caption{Principal implementations over the v12 logical format.}
  \label{tab:artifact}
  \begin{tabular}{@{}lll@{}}
    \toprule
    Component & Physical backend & Role \\
    \midrule
    Rust core & file or memory & format, MVCC, crypto, recovery \\
    Linux kernel VFS & raw block device & native mounted filesystem \\
    FUSE adapter & container file on host FS & portable mounted filesystem \\
    Sync/SaaS & opaque network objects & hosted or peer replication \\
    WebAssembly adapter & in-memory image & browser/application VFS API \\
    C ABI/tools & Rust core & embedding, inspection, repair \\
    \bottomrule
  \end{tabular}
\end{table}

\subsection{Rust Format Authority}

The core engine forbids unsafe Rust and does not depend on a general-purpose
serialization framework.  Header, catalog, unit-record, writer-set, key-grant,
WAL, and eviction formats are length-checked binary codecs.  The file backend
uses positioned reads and writes.  The memory backend exposes the same byte
image through an in-memory vector.  Both follow the same open path, including
header authentication, catalog and allocator reconstruction, and WAL recovery.

The current reader accepts exactly format v12.  This clean-cut policy prevents
silent interpretation of development-era formats, but it is not a format
stability promise: no migration implementation is included.  Cross-language
constants and derivations, including the fragment-size staircase and cipher
contexts, are recorded as golden vectors.  The artifact also includes hostile
decoder tests and six coverage-guided fuzz targets for unit records, commits,
eviction records, writer sets, key grants, and NoSQL records.

\subsection{Raw-Device Linux VFS}

The kernel path is an independent C implementation mounted directly on a block
device.  It parses and publishes the same headers, catalogs, unit records,
fragment geometry, encryption contexts, and writer state as the Rust engine.
The driver integrates with the Linux page cache and provides the broad VFS
subset needed by the evaluated filesystem: files and directories, read/write,
truncate, rename, symlinks, attributes, extended attributes and ACLs, and
filesystem statistics.  The implementation is intentionally described as an
experimental VFS rather than as complete POSIX conformance.  Remaining alias and
link-count semantics are listed in \cref{sec:discussion}.

Keeping the driver independent is useful evidence for format specificity: the
kernel does not call into the Rust engine or mount a userspace proxy.  C/Rust
round-trip fixtures and byte-level known-answer vectors guard the shared
format.  Loaded-module, target-kernel, and KASAN campaigns remain external
artifact gates because hosted CI can compile but cannot safely load the module
or provide a disposable block device.

\subsection{Portable Host-File VFS}

The FUSE adapter embeds the Rust engine and stores one container file on a host
filesystem.  This adds a userspace transition and a second filesystem layer,
but works without a kernel module and carries the complete engine feature set.
The adapter coalesces writes into explicit publication batches and implements
sequence-detecting read-ahead: a handle prefetches only after observing
sequential access, avoiding the prior random-read overfetch and memory pressure.
The same adapter layer contains macFUSE and WinFsp integration points, although
this paper evaluates only the Linux FUSE path and does not claim real-platform
validation for macOS or Windows.

The mounted surface exposes conflicts through extended attributes and maps
engine integrity failures to I/O errors rather than returning fabricated data.
FUSE and the kernel driver must not mount the same image concurrently: the
userspace advisory lock is taken on the container-file inode and cannot
coordinate with a block-device mount.

\subsection{Encrypted Hosted Storage and Peer Service}

The synchronization crate contains the transport-independent reconciliation
engine.  Its in-memory transport is used for model and integration tests.  The
production network client uses a compact binary wire format over TLS.  Batched
fragment requests amortize HTTP and authentication overhead.  The hosted
service supports direct TLS over HTTP/1.1, HTTP/2, and HTTP/3, or operation
behind a trusted TLS-terminating proxy.  Authentication uses SRP-6a followed by
scoped bearer tokens.  No password-equivalent value is sent during login.

Persistent server state is itself stored through an \sfs{} engine for account
and frontier metadata, with a separate append-oriented block-blob store for
large opaque payloads.  This is server-side operational encryption in addition
to, not a replacement for, the client's end-to-end container encryption.  The
service can verify signed record projections when writer enforcement is
enabled.  The deployed binary enables that policy by default.  It still cannot
interpret client paths, VFS metadata, or content in the secure profile.

The peer daemon reuses the hosted \texttt{/v1} wire and network client but
replaces the server store with an \texttt{EngineTransport} over a live
container.  A root-key-possession challenge replaces SRP.  This implementation
choice is what makes hosted and peer synchronization two topologies of one
protocol rather than two replication systems.

\subsection{WebAssembly and Embedded VFS}

The WebAssembly crate exposes container creation, password and raw-key opening,
file and directory creation, offset writes, truncate, read, list, and snapshot
operations through \texttt{wasm-bindgen}.  It uses the core memory backend, so
the byte vector returned by \texttt{snapshot} is the same container image that
the file backend and native VFS implementations open.  Password salts and
nonces route through the browser randomness backend.  The core disables its
native thread-based decrypt pool for wasm and uses the serial path.

The implementation establishes a source-level browser boundary, not yet a
browser deployment result.  Native tests exercise the wasm-selected core logic,
but the RC1 evidence does not include a hosted \texttt{wasm32-unknown-unknown}
build, JavaScript end-to-end test, or persistent OPFS/IndexedDB backend.  We
therefore call the bindings experimental throughout the paper.  Their relevance
is architectural: an application can own a complete encrypted VFS image and
hand only encrypted synchronization objects to the service, using the same
format and key model as the native clients.

\subsection{Artifact Evidence}

The source tree contains unit, property, integration, network, crash-seam, and
cross-format tests.  Relevant fixtures cover crashes before header publication
and round trips for GCM, XTS, and plaintext.  Further scenarios exercise signed
record and writer-set rejection, concurrent frontiers, third-replica
convergence, hosted and peer transports, and C/Rust golden vectors.  These tests
are evidence for implemented paths, not a formal proof of correctness or an
independent security audit.
